% ============================================================================
%  第 1 章 引言 —— body/01-introduction.tex
%  职责：讲清「为什么必须做」与「本文做了什么」，行文以问题驱动。
% ============================================================================

\section{引言}
\label{sec:introduction}

% ---------------------------------------------------------------------------
\subsection{研究背景}

无人机（\UAV）在城市低空的物流配送、设施巡检与应急响应中的应用规模持续增长，
随之而来的是其对地面人群、车辆与建筑物构成的 \termTPR。

\todo{补一到两句宏观数据：低空经济规模、飞行架次增速，注意引用权威统计源。}

% ---------------------------------------------------------------------------
\subsection{问题提出}

既有空域审批与航线规划多基于飞行前的静态快照，而实际运行环境随
时间尺度剧烈变化：

\begin{itemize}
  \item \textbf{微气象时变}：风速、风向与降雨强度改变失效概率；
  \item \textbf{建筑峡谷扰动}：街谷内的湍流与遮挡放大局部不确定性；
  \item \textbf{人口潮汐}：同一位置在通勤高峰与深夜的暴露人群差异达数量级。
\end{itemize}

由此形成「静态审批—动态环境」之间的安全裂隙，构成本文的核心问题：
\textit{如何在风险随时空变化的条件下，为给定出发时刻规划一条长期最优且在线可行的路径？}

% ---------------------------------------------------------------------------
\subsection{现有研究的局限}

从三个维度评述既有工作的不足（详见 \refsec{sec:related-work}）：
模型维度单一、时间维度缺失、求解维度保守。

% ---------------------------------------------------------------------------
\subsection{主要贡献}

本文的贡献可概括为四点：

\begin{enumerate}
  \item \textbf{模型}：构建……
  \item \textbf{算法}：提出 \methodnamelatex，……
  \item \textbf{理论}：证明……
  \item \textbf{实证}：在……场景验证……
\end{enumerate}

\todo{贡献清单要与 \refsec{sec:related-work} 的表格 Columns 一栏严丝合缝地对应。}

% ---------------------------------------------------------------------------
\subsection{文章结构}

全文组织如下：\refsec{sec:related-work} 回顾相关工作；
\refsec{sec:preliminaries} 给出问题定义；
\refsec{sec:risk-model} 与 \refsec{sec:planning} 分别阐述模型与算法；
\refsec{sec:experiments} 给出实验；\refsec{sec:discussion} 讨论；
\refsec{sec:conclusion} 总结。
